\section{Proposed Work}

For this thesis I propose to build a \textbf{verified compiler framework} for quantum programs grounded in the \textbf{operational semantics of \texttt{qstack}}.  
The goal is to formally validate the correctness of both \textbf{instruction-level} and \textbf{whole-circuit} compiler passes without relying on interactive theorem provers.  
Instead, the project will use a combination of \textbf{stabilizer-based logical-action verification} and \textbf{density-matrix equivalence checking} to automatically establish semantic preservation.  
For scalability, we prefer the former, so we will develop a tool to determine whether a circuit belongs to the stabilizer formalism.

\subsection{Language Design}

We will define a \textbf{minimal quantum language} derived from OpenQASM~3 that preserves the same \textbf{quantum/classical separation} as \texttt{qstack} while removing features that complicate formal reasoning, such as dynamic allocation and timing control.  
The retained constructs include:
\begin{itemize}
    \item A single universal \textbf{base gate};
    \item \textbf{User-defined composite gates} built from existing ones;
    \item \textbf{Static classical and quantum registers};
    \item \textbf{Gate application}, \textbf{measurement}, and \textbf{reset};
    \item \textbf{Conditionals} based on XORs of classical bits; and
    \item \textbf{\texttt{extern}} declarations for opaque operations.
\end{itemize}

This subset remains expressive enough to encode all \texttt{qstack} kernels and the \textbf{stabilizer circuits} defined in~\cite{kliuchnikov2023stabilizer}.  
It provides a compact foundation for formal semantics and automated verification.

\subsection{Formal Semantics}

The formal semantics will be inspired on the \textbf{small-step operational semantics} of \texttt{qstack}.  
A program configuration consists of:
\begin{itemize}
    \item A \textbf{quantum state} represented as a density matrix $\rho$;
    \item A \textbf{mapping} from logical qubit identifiers to physical indices; and
    \item A \textbf{classical store} containing deterministic registers and callback state.
\end{itemize}

Gate applications correspond to \textbf{unitary superoperators}, measurements collapse the relevant qubit and update the store, and callbacks may determine continuation through classical control.  
This semantics supports reasoning about both low-level physical transformations and high-level logical behavior under error correction.

As in \texttt{qstack}, evaluation of \texttt{extern} functions will be treated as opaque: the compiler assumes they return a list of quantum instructions but does not attempt to verify them.  
The output of callbacks will not be verified as part of the compilation process; I expect this is not a blocker since they should typically be used to handle only error-correction feedback, not core program semantics.

\subsection{Compiler Framework}

The compiler will be structured around two categories of passes that differ in both scope and verification method. Local, instruction-level passes will be statically verified once, while global, whole-circuit passes will rely on \emph{translation validation}.

\paragraph{Instruction-level passes.}
These are \emph{local transformations} that operate on one instruction at a time.  
Each pass defines a rule that replaces a specific instruction pattern with an equivalent one.  
For example, a canonicalization pass may replace every occurrence of \(H\) with the sequence \(S\!-\!SX\!-\!S\), leaving all other gates unchanged.  
These transformations are defined syntactically over single gates and are designed to preserve semantics in all contexts.  
Instruction-level passes include canonicalization, gate decomposition, and introduction of encoded (logical) instructions.  
Because each rule acts locally, these passes compose naturally to form larger compilation pipelines.

\paragraph{Whole-circuit passes.}
These are \emph{global transformations} that operate on an entire circuit at once.  
They may restructure or synthesize large regions of code.  
Typical examples include circuit optimization, scheduling, and layout mapping.  
Since these passes transform the full circuit in a single step, they are verified through direct equivalence checking between the input and output circuits rather than by composition of local proofs.

\subsection{Verification Methodology}

The verification methodology follows directly from this division between local and global transformations.
Instruction-level transformations admit static proofs of correctness because their effects are local and context-independent. In contrast, whole-circuit passes may restructure programs globally, making static proofs impractical; these are validated per compilation run using \textbf{translation validation}, which checks that the transformed circuit is equivalent to its input.

\paragraph{Instruction-level passes.}
Each instruction-level pass is verified in two stages:

\begin{enumerate}
  \item \textbf{Local equivalence proof.}  
  For the transformation under consideration (e.g., \(H \rightarrow S\!-\!SX\!-\!S\)), I'll prove that both sides have identical denotations using  \textbf{stabilizer verification} to show both circuits produce the same tableau transformation on the Pauli group.

  \item \textbf{Lifting to global correctness.}  
  Once local equivalence is established, a \textbf{forward-simulation argument} based on the operational semantics is used to prove that applying this transformation to every matching instance throughout a program preserves the program’s overall semantics.  
  Since each step preserves semantics, their composition preserves correctness for the entire compiled circuit.
\end{enumerate}

These passes are verified once, statically, as part of the compiler definition.

\paragraph{Whole-circuit passes.}
For transformations that rewrite or synthesize an entire circuit, verification is performed by direct \textbf{circuit equivalence checking} similarly by applying \textbf{stabilizer verification} when both input and output are Stabilizer circuits, by comparing their induced logical actions on the Pauli group

Because these transformations are verified as a single step, forward simulation is not required.  
Instead, we will use translation validation to check the equivalence of the full input and output circuits. 

\vspace{2em}

The focus of this thesis will be verification of stabilizer circuits. Verification of non-stabilizer circuits is possible by extending this framework using density-matrix or superoperator equivalence checking. However, this component will serve as a stretch goal, exploring how the same verification principles can generalize beyond stabilizer-based reasoning.
